% ============================================================
%  Preliminary Report: Diabetic Retinopathy Classification
%  Template -- fill in TODO sections before submission
% ============================================================
\documentclass[12pt, a4paper]{article}

% --- Packages -----------------------------------------------
\usepackage[utf8]{inputenc}
\usepackage[T1]{fontenc}
\usepackage[margin=1in]{geometry}
\usepackage{graphicx}
\usepackage{subcaption}
\usepackage{amsmath, amssymb}
\usepackage{booktabs}
\usepackage{multirow}
\usepackage{array}
\usepackage{float}
\usepackage{xcolor}
\usepackage{setspace}
\usepackage{enumitem}
\usepackage{hyperref}
\usepackage[capitalise]{cleveref}

\onehalfspacing

\hypersetup{
    colorlinks=true,
    linkcolor=blue,
    citecolor=blue,
    urlcolor=blue
}

% Placeholder box for missing figures
\newcommand{\figplaceholder}[2][6cm]{%
    \fbox{\parbox[c][#1][c]{\linewidth}{%
        \centering\textcolor{gray}{\small\textit{[Figure placeholder: #2]}}%
    }}%
}

% --- Title block --------------------------------------------
\title{
    \Large\textbf{Automated Grading of Diabetic Retinopathy\\
    Using Deep Convolutional Neural Networks:\\
    A Preliminary Study}\\[1em]
    \normalsize Preliminary Report
}
\author{
    Muhammad Uzair Khan\\
    \small\textit{TODO: Department / Institution}\\
    \small\texttt{mukskhan9999@gmail.com}
}
\date{\today}

% ============================================================
\begin{document}
\maketitle
\tableofcontents
\newpage

% ============================================================
\begin{abstract}
Diabetic retinopathy (DR) is a leading cause of preventable blindness worldwide.
Early and accurate grading of DR severity from fundus photographs is critical for
timely clinical intervention but remains resource-intensive and subject to
inter-grader variability.
This report presents a preliminary investigation into automated DR grading using
deep convolutional neural networks (CNNs).
We evaluate three backbone architectures --- ResNet-18, ResNet-50, and EfficientNet-B0
--- pretrained on ImageNet and fine-tuned on a fundus image dataset comprising five
severity grades (Grade 0--4).
We also study the effect of class-imbalance mitigation via weighted cross-entropy
loss and the utility of grayscale preprocessing.
Our best baseline model (EfficientNet-B0) achieves a Quadratic Weighted Kappa (QWK)
of 0.8771 on the held-out validation set.
Grad-CAM visualisations confirm that the models attend to clinically relevant
lesion regions including microaneurysms, haemorrhages, and neovascularisation.
\end{abstract}

% ============================================================
\section{Introduction}
\label{sec:intro}

Diabetes mellitus is a global epidemic affecting over 400 million people, with
diabetic retinopathy among its most prevalent microvascular complications.
DR manifests as progressive damage to the retinal vasculature and, if untreated,
leads to irreversible vision loss.
The International Clinical DR Disease Severity Scale classifies the condition into
five ordinal grades: No DR, Mild, Moderate, Severe non-proliferative DR (NPDR), and
Proliferative DR (PDR).

Manual grading by trained ophthalmologists is time-consuming and suffers from
inter-rater disagreement, especially in borderline cases.
Automated computer-aided diagnosis (CAD) systems based on deep learning have emerged
as a promising complement to clinical workflows, enabling high-throughput screening
in resource-limited settings.

\subsection{Problem Statement}
\label{ssec:problem}

Given a colour fundus photograph, the task is to predict the DR severity grade
$y \in \{0, 1, 2, 3, 4\}$.
The five-class ordinal nature of the labels --- where misclassifying Grade~1 as
Grade~3 is more severe than misclassifying it as Grade~2 --- motivates the use of
the Quadratic Weighted Kappa (QWK) as the primary evaluation metric, consistent
with the Kaggle DR Detection challenge \cite{TODO_kaggle}.

\subsection{Objectives}
\label{ssec:objectives}

\begin{enumerate}[label=\arabic*.]
    \item Establish strong CNN baselines under a unified training protocol.
    \item Investigate the impact of model capacity (ResNet-18 vs.\ ResNet-50 vs.\
          EfficientNet-B0) on ordinal classification performance.
    \item Assess the effect of class-imbalance mitigation strategies.
    \item Examine the role of colour versus grayscale input representations.
    \item Provide qualitative interpretability via Grad-CAM visualisations.
\end{enumerate}

\subsection{Report Structure}
The remainder of this report is organised as follows.
\Cref{sec:litreview} reviews related work.
\Cref{sec:dataset} describes the dataset.
\Cref{sec:methodology} details the experimental methodology.
\Cref{sec:results} presents and analyses results.
\Cref{sec:discussion} discusses findings and limitations.
\Cref{sec:conclusion} concludes with directions for future work.

% ============================================================
\section{Literature Review}
\label{sec:litreview}

\subsection{Early Automated DR Detection}
\label{ssec:lit_early}

Early automated DR detection relied on hand-crafted features extracted from the
optic disc, blood vessels, hard exudates, and microaneurysms.
TODO Author et al.\ \cite{TODO_early1} proposed \ldots
% TODO: summarise 1-2 classical machine-learning papers

\subsection{Deep Learning for Fundus Image Analysis}
\label{ssec:lit_dl}

The landmark work of Gulshan et al.\ \cite{TODO_gulshan2016} demonstrated that a
deep CNN trained on over 128,000 fundus images could achieve ophthalmologist-level
sensitivity and specificity for diabetic retinopathy.
Subsequent studies confirmed that transfer learning from large natural-image
datasets provides a strong initialisation for fundus analysis
\cite{TODO_transfer_learning}.

\subsubsection{Backbone Architectures}
ResNet \cite{TODO_he2016} introduced residual connections that allow training of
very deep networks by mitigating the vanishing-gradient problem.
EfficientNet \cite{TODO_tan2019} proposed a principled compound scaling strategy
that simultaneously scales network depth, width, and resolution, achieving
superior accuracy-efficiency trade-offs.
% TODO: cite additional architecture papers relevant to DR.

\subsubsection{Class Imbalance in DR Datasets}
DR datasets are inherently imbalanced --- a majority of patients exhibit no or mild
retinopathy, whereas advanced grades are rare.
Strategies reported in the literature include weighted loss functions
\cite{TODO_weighted_loss}, oversampling (SMOTE, Mixup), and data augmentation
tailored to medical images \cite{TODO_augmentation}.

\subsubsection{Ordinal Classification}
Because DR grades form an ordinal scale, standard cross-entropy loss does not
exploit the ordering information.
CORN \cite{TODO_corn} and CORAL \cite{TODO_coral} reformulate ordinal regression
as a set of binary sub-tasks, enforcing rank consistency in the output layer.
Focal loss \cite{TODO_focal} down-weights easy examples and has shown benefits for
imbalanced medical imaging tasks.

\subsection{Interpretability in Medical AI}
\label{ssec:lit_xai}

Class Activation Mapping (CAM) and its gradient-based extension, Grad-CAM
\cite{TODO_gradcam}, produce saliency maps that highlight image regions most
influential to a model's prediction.
In ophthalmology, Grad-CAM has been used to verify that CNNs attend to known
lesion types rather than image artefacts \cite{TODO_gradcam_eye}.

\subsection{Summary and Research Gap}
\label{ssec:lit_gap}

% TODO: 2-3 sentences identifying what prior work leaves open and how this work addresses it
Existing studies have \ldots
This work contributes \ldots

% ============================================================
\section{Dataset}
\label{sec:dataset}

\subsection{Source and Description}
\label{ssec:data_source}

% TODO: replace with the actual dataset name/link/citation
The dataset used in this study consists of colour fundus photographs labelled by
trained graders into five DR severity grades.
\Cref{tab:class_dist} summarises the per-class sample counts.

\begin{table}[H]
\centering
\caption{Class distribution in the dataset.}
\label{tab:class_dist}
\begin{tabular}{lcc}
\toprule
\textbf{Grade} & \textbf{Label} & \textbf{No.\ Images} \\
\midrule
0 & No DR          & 1805 \\
1 & Mild NPDR      & TODO \\
2 & Moderate NPDR  & TODO \\
3 & Severe NPDR    & 193  \\
4 & Proliferative  & TODO \\
\midrule
  & \textbf{Total} & TODO \\
\bottomrule
\end{tabular}
\end{table}

\subsection{Preprocessing}
\label{ssec:preprocessing}

All images were resized to $224 \times 224$ pixels using bicubic interpolation.
Pixel values were normalised using ImageNet channel statistics
($\mu = [0.485, 0.456, 0.406]$, $\sigma = [0.229, 0.224, 0.225]$).
For the grayscale experiment (Exp~5), images were converted to single-channel
luminance representations and the channel mean/std were adjusted accordingly.

\subsection{Data Augmentation}
\label{ssec:augmentation}

The following augmentations were applied on-the-fly during training:
random horizontal and vertical flipping, random rotation ($\pm 30^\circ$),
and colour jitter (brightness, contrast, saturation, hue).
No augmentation was applied at inference time.

\subsection{Train / Validation Split}
\label{ssec:split}

The dataset was split into 80\% training and 20\% validation using stratified
sampling to preserve class proportions across both subsets.

% ============================================================
\section{Methodology}
\label{sec:methodology}

\subsection{Experimental Protocol}
\label{ssec:protocol}

All experiments share the following hyperparameter configuration unless stated
otherwise:
\begin{itemize}
    \item Optimiser: Adam ($\beta_1=0.9$, $\beta_2=0.999$, $\epsilon=10^{-8}$)
    \item Learning rate: $10^{-4}$ (no scheduler)
    \item Batch size: 16 (Exp~1--4); 8 (Exp~5)
    \item Training epochs: 10
    \item Hardware: TODO (GPU type, VRAM)
\end{itemize}

\subsection{Baseline --- ResNet-18 (Exp~1)}
\label{ssec:exp1}

ResNet-18 \cite{TODO_he2016} (11M parameters) pretrained on ImageNet serves as
the primary baseline.
The final fully connected layer was replaced with a linear layer mapping to 5
classes.
Standard cross-entropy loss was used.

\subsection{ResNet-50 (Exp~2)}
\label{ssec:exp2}

ResNet-50 (25M parameters) was evaluated to assess the effect of increased model
capacity while keeping the architecture family constant.

\subsection{EfficientNet-B0 (Exp~3)}
\label{ssec:exp3}

EfficientNet-B0 (5.3M parameters) was included to test whether compound-scaled
architectures offer a better accuracy-to-parameter trade-off on this task.

\subsection{Weighted Cross-Entropy (Exp~4)}
\label{ssec:exp4}

To address class imbalance, per-class weights were computed as the inverse of
normalised class frequencies:
\begin{equation}
    w_c = \frac{N}{\,K \cdot n_c\,},
\end{equation}
where $N$ is the total number of training samples, $K$ is the number of classes,
and $n_c$ is the count of class $c$.
This experiment used the ResNet-18 backbone.

\subsection{Grayscale Preprocessing (Exp~5)}
\label{ssec:exp5}

To investigate whether colour information is essential for DR grading, images were
converted to grayscale and the ResNet-18 input layer was adapted to accept a
single channel.
This also serves as a domain-alignment ablation, since many clinical fundus cameras
produce green-channel or grayscale images.

\subsection{Evaluation Metrics}
\label{ssec:metrics}

\paragraph{Quadratic Weighted Kappa (QWK).}
The primary metric, sensitive to the ordinal structure of DR grades:
\begin{equation}
    \kappa = 1 - \frac{\sum_{i,j} w_{ij} O_{ij}}{\sum_{i,j} w_{ij} E_{ij}},
    \quad w_{ij} = \frac{(i-j)^2}{(K-1)^2}.
\end{equation}
A QWK of 0 indicates chance agreement; 1 indicates perfect agreement.

\paragraph{Macro-averaged F1-score.}
Complements QWK by treating each class equally regardless of support.

\paragraph{Per-class precision, recall, and F1.}
Reported in the classification report to diagnose per-grade performance.

\paragraph{Confusion matrix.}
Visualised to reveal systematic grade confusion patterns.

\subsection{Interpretability: Grad-CAM}
\label{ssec:gradcam}

Gradient-weighted Class Activation Maps (Grad-CAM) \cite{TODO_gradcam} were
computed for the final convolutional layer of each backbone to produce heatmaps
highlighting discriminative image regions for each predicted grade.

% ============================================================
\section{Results}
\label{sec:results}

\subsection{Quantitative Comparison}
\label{ssec:quant}

\Cref{tab:main_results} summarises the performance of all experiments on the
validation set.

\begin{table}[H]
\centering
\caption{Validation performance across experiments. Best values in \textbf{bold}.}
\label{tab:main_results}
\begin{tabular}{llccc}
\toprule
\textbf{Exp} & \textbf{Model / Setting} & \textbf{QWK} & \textbf{Accuracy (\%)} & \textbf{Macro F1} \\
\midrule
1 & ResNet-18 (baseline)      & TODO            & TODO & TODO \\
2 & ResNet-50                 & 0.8678          & TODO & TODO \\
3 & EfficientNet-B0           & \textbf{0.8771} & TODO & TODO \\
4 & ResNet-18 + Weighted CE   & 0.8666          & TODO & TODO \\
5 & ResNet-18 + Grayscale     & 0.8690          & TODO & TODO \\
\bottomrule
\end{tabular}
\end{table}

\subsection{Training Dynamics}
\label{ssec:training_curves}

\begin{figure}[H]
    \centering
    \begin{subfigure}[b]{0.48\textwidth}
        \centering
        \figplaceholder{training\_curves\_resnet18.png}
        \caption{ResNet-18 (Exp~1)}
    \end{subfigure}
    \hfill
    \begin{subfigure}[b]{0.48\textwidth}
        \centering
        \figplaceholder{training\_curves\_efficientnet.png}
        \caption{EfficientNet-B0 (Exp~3)}
    \end{subfigure}
    \caption{Training and validation loss / QWK curves for selected experiments.}
    \label{fig:training_curves}
\end{figure}

\subsection{Confusion Matrices}
\label{ssec:conf_matrices}

\begin{figure}[H]
    \centering
    \begin{subfigure}[b]{0.32\textwidth}
        \centering
        \figplaceholder[5cm]{confusion\_matrix\_resnet18.png}
        \caption{ResNet-18}
    \end{subfigure}
    \hfill
    \begin{subfigure}[b]{0.32\textwidth}
        \centering
        \figplaceholder[5cm]{confusion\_matrix\_efficientnet.png}
        \caption{EfficientNet-B0}
    \end{subfigure}
    \hfill
    \begin{subfigure}[b]{0.32\textwidth}
        \centering
        \figplaceholder[5cm]{confusion\_matrix\_weighted\_ce.png}
        \caption{ResNet-18 + Weighted CE}
    \end{subfigure}
    \caption{Normalised confusion matrices (rows: true grade; columns: predicted grade).}
    \label{fig:conf_matrices}
\end{figure}

\subsection{Grad-CAM Visualisations}
\label{ssec:gradcam_results}

\begin{figure}[H]
    \centering
    \figplaceholder[7cm]{gradcam\_examples\_resnet18.png}
    \caption{Grad-CAM overlays for ResNet-18 across DR grades 0--4.
             Warmer colours indicate regions of higher importance to the prediction.}
    \label{fig:gradcam}
\end{figure}

% ============================================================
\section{Discussion}
\label{sec:discussion}

\subsection{Effect of Architecture}
\label{ssec:disc_arch}

EfficientNet-B0 outperformed both ResNet variants despite having fewer parameters
than ResNet-18, suggesting that compound scaling leads to more efficient feature
extraction for retinal images.
The marginal gap between ResNet-50 and EfficientNet-B0 (TODO pp.\ QWK) implies
that \ldots
% TODO: expand analysis

\subsection{Effect of Class-Imbalance Weighting}
\label{ssec:disc_imbalance}

Weighted cross-entropy improved \ldots / did not substantially improve \ldots
% TODO: discuss per-class recall changes, especially for minority Grades 3 and 4

\subsection{Role of Colour Information}
\label{ssec:disc_colour}

The grayscale experiment achieved competitive performance (QWK 0.8690), indicating
that \ldots
% TODO: discuss clinical implications -- haemorrhage detection in red channel, etc.

\subsection{Interpretability Analysis}
\label{ssec:disc_xai}

Grad-CAM maps confirm that the model attends to clinically meaningful structures:
microaneurysms and dot haemorrhages in Grade~1--2 images, large haemorrhages and
venous beading in Grade~3, and neovascularisation near the optic disc in Grade~4.
% TODO: elaborate with specific examples from the figure

\subsection{Limitations}
\label{ssec:limitations}

\begin{itemize}
    \item The dataset is highly imbalanced; Grade~3 contains only 193 images,
          limiting reliable evaluation of minority-class performance.
    \item Training was limited to 10 epochs without a learning-rate schedule;
          longer training with cosine annealing or cyclical LR may improve results.
    \item No test-time augmentation (TTA) was applied.
    \item All models were evaluated on a single random split; cross-validation
          would yield more reliable estimates.
    \item TODO: add any other limitations specific to your setup.
\end{itemize}

% ============================================================
\section{Conclusion and Future Work}
\label{sec:conclusion}

This preliminary study established five CNN baselines for five-class diabetic
retinopathy grading from colour fundus photographs.
EfficientNet-B0 achieved the best QWK of 0.8771, outperforming deeper ResNet
variants while using fewer parameters.
Weighted cross-entropy and grayscale preprocessing each offered modest gains over
the ResNet-18 baseline.
Grad-CAM visualisations confirmed that the learned representations are clinically
interpretable.

\subsection*{Future Work}

\begin{enumerate}[label=\arabic*.]
    \item \textbf{Ordinal loss functions} (CORAL, CORN) to exploit label ordering.
    \item \textbf{Focal loss} to further mitigate class imbalance.
    \item \textbf{Advanced data augmentation}: Mixup, CutMix, or domain-specific
          green-channel enhancement.
    \item \textbf{Stronger backbones}: EfficientNet-B4/B7, ConvNeXt, or
          Vision Transformers (ViT, Swin).
    \item \textbf{Learning-rate scheduling} and longer training runs.
    \item \textbf{Ensemble} top-performing models to further boost QWK.
    \item \textbf{External validation} on an independent dataset to assess
          generalisation.
\end{enumerate}

% ============================================================
\begin{thebibliography}{99}

\bibitem{TODO_kaggle}
Kaggle, ``Diabetic Retinopathy Detection Challenge,'' 2015.
[Online]. Available: \url{https://www.kaggle.com/c/diabetic-retinopathy-detection}

\bibitem{TODO_early1}
TODO Author et al., ``TODO: Title of classical DR detection paper,''
\textit{TODO Journal / Conference}, vol.~TODO, pp.~TODO--TODO, TODO year.

\bibitem{TODO_gulshan2016}
V.~Gulshan et al., ``Development and Validation of a Deep Learning Algorithm
for Detection of Diabetic Retinopathy in Retinal Fundus Photographs,''
\textit{JAMA}, vol.~316, no.~22, pp.~2402--2410, 2016.

\bibitem{TODO_transfer_learning}
TODO Author et al., ``TODO: Transfer learning for fundus image analysis,''
\textit{TODO Journal / Conference}, TODO year.

\bibitem{TODO_he2016}
K.~He, X.~Zhang, S.~Ren, and J.~Sun, ``Deep Residual Learning for Image
Recognition,'' in \textit{Proc.\ CVPR}, Las Vegas, NV, 2016, pp.~770--778.

\bibitem{TODO_tan2019}
M.~Tan and Q.~Le, ``EfficientNet: Rethinking Model Scaling for Convolutional
Neural Networks,'' in \textit{Proc.\ ICML}, 2019.

\bibitem{TODO_weighted_loss}
TODO Author et al., ``TODO: Weighted loss / class imbalance in medical imaging,''
\textit{TODO Journal / Conference}, TODO year.

\bibitem{TODO_augmentation}
TODO Author et al., ``TODO: Data augmentation strategies for retinal images,''
\textit{TODO Journal / Conference}, TODO year.

\bibitem{TODO_corn}
TODO Author et al., ``TODO: CORN --- Conditional Ordinal Regression for NNs,''
\textit{TODO Journal / Conference}, TODO year.

\bibitem{TODO_coral}
W.~Cao, V.~Mir\'{e}i, J.~Ramirez-de-Arellano, and S.~Raschka,
``Rank Consistent Ordinal Regression for Neural Networks with Application to
Age Estimation,'' \textit{Pattern Recognition Letters}, 2020.

\bibitem{TODO_focal}
T.-Y.~Lin, P.~Goyal, R.~Girshick, K.~He, and P.~Doll\'{a}r,
``Focal Loss for Dense Object Detection,''
\textit{IEEE Trans.\ Pattern Anal.\ Mach.\ Intell.}, vol.~42, no.~2, 2020.

\bibitem{TODO_gradcam}
R.~R.~Selvaraju et al., ``Grad-CAM: Visual Explanations from Deep Networks
via Gradient-based Localization,'' \textit{Int.\ J.\ Comput.\ Vis.}, vol.~128,
pp.~336--359, 2020.

\bibitem{TODO_gradcam_eye}
TODO Author et al., ``TODO: Grad-CAM applied to ophthalmology / retinal imaging,''
\textit{TODO Journal / Conference}, TODO year.

\end{thebibliography}

\end{document}
